\section{Experiments}

The experiments are designed as audit and boundary diagnostics, not as a
benchmark-superiority study.

\paragraph{Hidden posterior conflict.}
WNUT17 diagnostic stress tests show that legal constrained decoding can coexist
with low posterior BIO mass.  These runs are diagnostic only; they do not
support a NER F1 improvement claim.

% Generated by scripts/analysis/generate_latex_tables.py.
% Do not edit numeric values by hand; edit curated CSVs and regenerate.
\begin{table}[t]
\centering
\small
\caption{WNUT17 diagnostic and viability checks.}
\label{tab:table:2:r5:wnut17}
\begin{tabular}{llllllll}
\toprule
block & variant & seeds & P\_event & delta\_P & hidden\_conflict & entity\_F1 & claim\_use \\
\midrule
r5a\_diagnostic\_stress & B0\_unconstrained & 10 & 0.0566 & 0.0000 & 1.0000 & 0.0000 & hidden-conflict diagnostic evidence \\
r5a\_diagnostic\_stress & B4\_semi\_event\_0.1 & 10 & 0.3389 & 0.2822 & 0.9963 & 0.0000 & hidden-conflict diagnostic evidence \\
r5a\_diagnostic\_stress & B5\_rule\_feature\_0.8 & 10 & 0.1608 & 0.1042 & 1.0000 & 0.0000 & hidden-conflict diagnostic evidence \\
r5a\_diagnostic\_stress & B6\_pr\_style\_tau0.7 & 10 & 0.1703 & 0.1137 & 1.0000 & 0.0000 & hidden-conflict diagnostic evidence \\
r5b\_feature\_viability & B0\_unconstrained & 10 & 0.9824 & 0.0000 & 0.0088 & 0.1660 & task viability boundary \\
r5b\_feature\_viability & B4\_semi\_event\_0.1 & 10 & 0.9865 & 0.0041 & 0.0074 & 0.1522 & task viability boundary \\
r5b\_feature\_viability & B5\_rule\_feature\_0.8 & 10 & 0.9864 & 0.0040 & 0.0058 & 0.1645 & task viability boundary \\
r5b\_feature\_viability & B6\_pr\_style\_tau0.7 & 10 & 0.9868 & 0.0044 & 0.0060 & 0.1463 & task viability boundary \\
\bottomrule
\end{tabular}
\end{table}


\paragraph{Event-mass movement.}
Controlled, semi-real, real-source, and public sequence-labeling cases show
that event training can raise posterior event mass.  Task metrics vary, so the
claim is event-mass movement rather than task superiority.

% Generated by scripts/analysis/generate_latex_tables.py.
% Do not edit numeric values by hand; edit curated CSVs and regenerate.
\begin{table}[t]
\centering
\small
\caption{Event-training movement in controlled and field-style probes.}
\label{tab:table:3:event:training:signal}
\begin{tabular}{llllllll}
\toprule
block & family & rows & mean\_delta\_P & positive\_P\_rate & mean\_delta\_legal & mean\_delta\_char & mean\_delta\_exact \\
\midrule
r1\_controlled\_formal & B4 & 8 & 0.1883 & 1 & 0.1639 & 0.0048 & 0.0032 \\
r2\_semi\_real\_formal & B4 & 8 & 0.1584 & 1 & 0.0854 & 0.0188 & 0.0775 \\
r2\_semi\_real\_formal & B5 & 4 & 0.0740 & 1 & 0.0671 & 0.0119 & 0.0449 \\
r2\_semi\_real\_formal & B6 & 4 & 0.0968 & 1 & 0.0772 & 0.0210 & 0.0723 \\
r4\_real\_source\_formal & B4 & 6 & 0.0722 & 1 & 0.0018 & 0.0289 & 0.1127 \\
r4\_real\_source\_formal & B5 & 3 & 0.0314 & 1 & 0.0013 & 0.0187 & 0.0678 \\
r4\_real\_source\_formal & B6 & 3 & 0.0219 & 1 & 0.0009 & 0.0190 & 0.0743 \\
\bottomrule
\end{tabular}
\end{table}


\paragraph{Diagnostic uncertainty boundary.}
In evaluated field-style diagnostics, event risk has positive ranking signal,
but generic uncertainty baselines are stronger in several settings.  We
therefore report event risk as a rule-specific audit statistic, not a universal
uncertainty replacement.

% Generated by scripts/analysis/generate_latex_tables.py.
% Do not edit numeric values by hand; edit curated CSVs and regenerate.
\begin{table}[t]
\centering
\small
\caption{R6a diagnostic ranking and uncertainty boundary.}
\label{tab:table:4:diagnostic:ranking}
\begin{tabular}{lllllllll}
\toprule
source & task & cases & base\_exact\_error & AUROC & AUPRC & Spearman & bottom20\_exact\_error & bottom20\_hidden\_conflict \\
\midrule
all & all & 21000 & 0.7359 & 0.7088 & 0.8470 & 0.2869 & 0.8586 & 0.1033 \\
real\_source & invoice\_6d & 3000 & 0.7203 & 0.7928 & 0.8862 & 0.4949 & 0.9150 & 0 \\
real\_source & invoice\_c6d & 3000 & 0.6547 & 0.7354 & 0.8116 & 0.4347 & 0.8333 & 0.0067 \\
real\_source & stock\_5d & 3000 & 0.7143 & 0.8334 & 0.8979 & 0.5667 & 0.9167 & 0 \\
semi\_real & amount & 3000 & 0.7810 & 0.8046 & 0.9220 & 0.5178 & 0.9600 & 0.0750 \\
semi\_real & date & 3000 & 0.5743 & 0.7220 & 0.7918 & 0.4001 & 0.8817 & 0.4783 \\
semi\_real & dose & 3000 & 0.7867 & 0.8509 & 0.9367 & 0.5296 & 0.9400 & 0.0067 \\
semi\_real & product\_code & 3000 & 0.9200 & 0.8240 & 0.9754 & 0.5165 & 0.9800 & 0.1567 \\
\bottomrule
\end{tabular}
\end{table}

% Generated by scripts/analysis/generate_latex_tables.py.
% Do not edit numeric values by hand; edit curated CSVs and regenerate.
\begin{table}[t]
\centering
\small
\caption{CoNLL2000 public uncertainty boundary.}
\label{tab:table:7:public:uncertainty}
\begin{tabular}{llllll}
\toprule
variant & score & AUROC\_exact & AUPRC\_exact & Spearman\_token\_error & exact\_risk\_gap \\
\midrule
B0\_unconstrained & event\_risk\_1\_minus\_p & 0.7384 & 0.8829 & 0.3117 & 0.4450 \\
B0\_unconstrained & token\_marginal\_entropy & 0.8305 & 0.9388 & 0.4559 & 0.5867 \\
B0\_unconstrained & sequence\_entropy & 0.8339 & 0.9405 & 0.4346 & 0.5817 \\
B0\_unconstrained & viterbi\_margin\_inverse & 0.7860 & 0.9040 & 0.3566 & 0.5167 \\
B0\_unconstrained & max\_sequence\_probability\_inverse & 0.8211 & 0.9334 & 0.4217 & 0.5717 \\
B0\_unconstrained & neg\_log\_viterbi\_probability & 0.8211 & 0.9334 & 0.4217 & 0.5717 \\
B4\_semi\_event\_0.1 & event\_risk\_1\_minus\_p & 0.6895 & 0.8865 & 0.2886 & 0.2650 \\
B4\_semi\_event\_0.1 & token\_marginal\_entropy & 0.8169 & 0.9382 & 0.3972 & 0.5300 \\
B4\_semi\_event\_0.1 & sequence\_entropy & 0.8178 & 0.9394 & 0.3785 & 0.5283 \\
B4\_semi\_event\_0.1 & viterbi\_margin\_inverse & 0.7749 & 0.9063 & 0.3072 & 0.4650 \\
B4\_semi\_event\_0.1 & max\_sequence\_probability\_inverse & 0.8030 & 0.9312 & 0.3590 & 0.5083 \\
B4\_semi\_event\_0.1 & neg\_log\_viterbi\_probability & 0.8030 & 0.9312 & 0.3590 & 0.5083 \\
\bottomrule
\end{tabular}
\end{table}


\paragraph{Public CoNLL2000 case study.}
The public BIO/chunking case provides a multiseed structured-prediction audit.
B4 raises mean posterior BIO event mass and slightly reduces hidden conflict,
while mean token accuracy and span F1 decrease.  This is boundary evidence
against reading event loss as a general accuracy method.

% Generated by scripts/analysis/generate_latex_tables.py.
% Do not edit numeric values by hand; edit curated CSVs and regenerate.
\begin{table}[t]
\centering
\small
\caption{CoNLL2000 public BIO/chunking multiseed case study.}
\label{tab:table:6:public:conll2000}
\begin{tabular}{llllllllll}
\toprule
setting & variant & seeds & P\_BIO\_mean & P\_BIO\_std & hidden\_conflict\_mean & B7\_legal\_mean & token\_acc\_mean & span\_F1\_mean & claim\_use \\
\midrule
formal\_3seed & B0\_unconstrained & 3 & 0.9837 & 0.0011 & 0.0133 & 1.0000 & 0.8731 & 0.7973 & public structured-prediction multiseed audit case; not benchmark superiority \\
formal\_3seed & B4\_semi\_event\_0.1 & 3 & 0.9885 & 0.0063 & 0.0123 & 1.0000 & 0.8697 & 0.7918 & public structured-prediction multiseed audit case; not benchmark superiority \\
\bottomrule
\end{tabular}
\end{table}


\paragraph{Constrained-product baseline and sensitivity.}
B7 reports a faithful constrained-product Viterbi baseline using the same BIO
language.  R7 sensitivity studies show that rule and lambda choice matter,
including a misleading-rule case where event mass rises while task metrics
fall.

% Generated by scripts/analysis/generate_latex_tables.py.
% Do not edit numeric values by hand; edit curated CSVs and regenerate.
\begin{table}[t]
\centering
\small
\caption{B7 constrained-product decoding baseline.}
\label{tab:table:8:b7:constrained:product}
\begin{tabular}{llllll}
\toprule
source\_model & decoded\_legal\_rate & token\_accuracy & entity\_F1 & uses\_event\_training & uses\_event\_mass\_for\_decoding \\
\midrule
B0\_unconstrained & 1.0000 & 0.8841 & 0.1486 & False & False \\
B4\_semi\_event\_0.1 & 1.0000 & 0.8838 & 0.1688 & True & False \\
\bottomrule
\end{tabular}
\end{table}

% Generated by scripts/analysis/generate_latex_tables.py.
% Do not edit numeric values by hand; edit curated CSVs and regenerate.
\begin{table}[t]
\centering
\small
\caption{R7 lambda/rule sensitivity and boundary cases.}
\label{tab:table:9:r7:sensitivity}
\begin{tabular}{lllllllll}
\toprule
task & variant & runs & P\_event & delta\_P\_event & delta\_legal\_rate & delta\_char\_acc & delta\_exact\_acc & boundary \\
\midrule
date & B0\_unconstrained & 10 & 0.7136 & 0.0000 & 0.0000 & 0.0000 & 0.0000 & ordinary \\
date & B4\_semi\_event\_1.0 & 10 & 0.9738 & 0.2603 & 0.1948 & 0.0306 & 0.1996 & ordinary \\
product\_code & B0\_unconstrained & 10 & 0.7765 & 0.0000 & 0.0000 & 0.0000 & 0.0000 & ordinary \\
product\_code & B4\_semi\_event\_1.0 & 10 & 0.9844 & 0.2079 & 0.0328 & 0.0097 & 0.0248 & ordinary \\
product\_code\_swapped\_rule & B0\_unconstrained & 10 & 0.0000 & 0.0000 & 0.0000 & 0.0000 & 0.0000 & event/task tradeoff \\
product\_code\_swapped\_rule & B4\_semi\_event\_0.1 & 10 & 0.6083 & 0.6083 & 0.6752 & -0.4352 & -0.0668 & event/task tradeoff \\
product\_code\_swapped\_rule & B4\_semi\_event\_1.0 & 10 & 0.9486 & 0.9485 & 0.9716 & -0.5524 & -0.0816 & event/task tradeoff \\
stock\_like\_digits & B0\_unconstrained & 10 & 0.9192 & 0.0000 & 0.0000 & 0.0000 & 0.0000 & legal-rate-not-useful \\
stock\_like\_digits & B4\_semi\_event\_1.0 & 10 & 0.9961 & 0.0769 & 0.0000 & 0.0331 & 0.0868 & legal-rate-not-useful \\
\bottomrule
\end{tabular}
\end{table}

